Este apêndice registra onde estão os dados usados para sustentar a análise
quantitativa do Capítulo~\ref{chapter:resultados-testes} e reúne os detalhes
operacionais da campanha experimental. A campanha \texttt{final} corresponde à
rodada consolidada de testes executada com a versão do \emph{firmware} analisada
nesta monografia. No repositório local, os artefatos estão em
\path{transmitter/testing/results/final}\footnote{Disponível em:
\url{https://github.com/cirillom/WCAN/tree/thesis-final/transmitter/testing/results/final}.}.

A pasta contém o plano de execução, a configuração das placas, a configuração das
suítes, o resumo produzido pelo \emph{runner}, o relatório consolidado do
analisador e os \emph{logs} seriais de cada placa. Cada execução possui um
\texttt{scenario.json} e um conjunto de \emph{logs} separados por papel e porta
serial. Essa organização permite recalcular as métricas agregadas e inspecionar
manualmente casos reprovados.

\section{Configuração e Plano de Execução}

O arquivo \texttt{boards.yaml} descreve o inventário físico da bancada, incluindo
porta serial, tipo de chip e identificador CAN inicial de cada placa. O arquivo
\texttt{tests.yaml} define as suítes, frequências, duração, repetições,
transportes, valor de \texttt{linger\_ms} e quantidade de identificadores CAN por
sensor. Antes da execução, o \emph{runner} gera \texttt{run\_plan.json}; para
cada ensaio, também é gerado um \texttt{scenario.json} com a topologia, os papéis
atribuídos às placas, os filtros esperados e os parâmetros enviados por UART.

Na campanha final, o arquivo \texttt{run\_plan.json} registra 684 execuções
planejadas, distribuídas entre as variantes \texttt{BROADCAST} e
\texttt{MULTICAST}, com duas repetições por cenário aplicável. O arquivo
\texttt{summary.csv} contém uma linha por execução da campanha, com 668 linhas
\texttt{OK} e 16 linhas \texttt{MONITOR\_ERROR}. O relatório
\texttt{analysis\_summary.txt} consolida 453 execuções aprovadas em 684 e 249
topologias aprovadas em 354 agrupamentos (330 com duas repetições e 24 de
execução única em \texttt{mixed\_frequency}), considerando aprovado o agrupamento
de topologia com ao menos uma repetição bem-sucedida. Para a interpretação do
protocolo no Capítulo~\ref{chapter:resultados-testes}, cinco reprovações por
corrupção de estatísticas ou \emph{log} serial foram separadas da base útil,
resultando em 453 aprovações sobre 679 execuções interpretáveis.

\section{Fluxo de Automação e Protocolo UART}

Os testes usam uma estrutura dedicada em \texttt{wcan\_test}, responsável por
configurar o papel de cada placa, iniciar a janela de medição, acionar o gerador
sintético de amostras e imprimir as métricas consumidas pelo analisador. A classe
\texttt{WcanTestSession} recebe um \texttt{TestConfig}, aguarda o comando
\texttt{WCAN\_TEST\_START} via UART, reinicia os contadores da camada de
estatísticas, inicia o \texttt{RampCanSensor} quando a placa atua como sensor,
mantém o ensaio ativo pelo tempo configurado e, ao final, para o sensor, drena o
transceptor, imprime as estatísticas e sinaliza \texttt{WCAN\_TEST\_END} para o
computador de testes. Como a sessão opera sobre uma referência para
\texttt{TransceiverBase}, o mesmo fluxo de teste é usado para as variantes
\emph{broadcast} e \emph{multicast}; a diferença entre elas fica encapsulada na
implementação do transceptor escolhida em tempo de compilação.

No \emph{firmware}, \texttt{main.cpp} atua como camada de inicialização. A função
\texttt{app\_main} aguarda a configuração de \emph{boot} via UART, inicializa NVS
e Wi-Fi quando a placa participa do teste e instancia o \texttt{Transceiver} de
acordo com o papel recebido. Em sensores, o transceptor é criado com a lista de
identificadores CAN de transmissão e com o \texttt{linger\_ms} do cenário. Em
receptores, o transceptor é criado com a lista de identificadores CAN aceitos
quando a filtragem ativa está habilitada. Placas ociosas apenas aguardam o
comando de início e encerram a execução, permitindo que a matriz use subconjuntos
da bancada sem alterar o \emph{firmware}.

A sincronização entre \emph{host} e placas é feita por um protocolo textual
simples em UART. Ao iniciar, cada placa imprime \texttt{WCAN\_CFG\_WAIT v=1}; o
\emph{runner} envia uma linha de configuração com papel, transporte, duração,
frequência, identificador CAN inicial, quantidade de identificadores CAN,
\texttt{linger\_ms} e filtros; o \emph{firmware} responde
\texttt{WCAN\_CFG\_ACK} quando a configuração é aceita. Depois que todas as
placas do cenário estão prontas, o \emph{host} envia \texttt{WCAN\_TEST\_START}.
Durante o teste, sensores usam \texttt{RampCanSensor}, que gera um contador
monotônico por \texttt{esp\_timer} periódico e chama \texttt{send\_data} para
cada identificador CAN configurado. Ao final da janela, o \emph{firmware} drena
o transceptor, imprime as métricas coletadas e encerra com
\texttt{WCAN\_TEST\_END}; abortos são registrados como
\texttt{WCAN\_TEST\_ABORT}.

\section{Estrutura dos Logs e Métricas}

A medição fica concentrada na hierarquia \texttt{Stats}. A classe base define os
pontos de instrumentação e, quando a \emph{build} é compilada com
\texttt{MEASURE}, \texttt{make\_stats()} retorna a implementação
\texttt{MeasuredStats}. Essa implementação registra eventos de recepção,
\emph{airtime}, agrupamento temporal, despacho de lotes e falhas de envio do
sensor. Para reduzir a pressão sobre a UART, os contadores recebidos e as falhas
são compactados em intervalos, como \texttt{R(RANGE)} e \texttt{S(FAIL)}, em vez
de imprimir uma linha para cada amostra individual. As estatísticas de lote e de
memória são emitidas ao final da execução por meio de \texttt{WCAN\_BATCH} e
\texttt{WCAN\_MEASURES}.

No lado sensor, a linha \texttt{WCAN\_SENSOR\_END} informa o último contador
gerado e a frequência média observada. No lado receptor, as linhas
\texttt{R(RANGE)} registram intervalos de contadores aceitos para cada
identificador CAN. O script \texttt{analysis.py} reconstrói os conjuntos de
contadores transmitidos e recebidos e compara, para cada sensor e identificador
CAN, o conjunto gerado com a união dos conjuntos recebidos pelos receptores
esperados.

Falhas observadas no caminho de transmissão também são registradas. As linhas
\texttt{P(FAIL)} indicam falhas de entrega de pacotes no mecanismo de transporte,
com destaque para os casos em que o \emph{broadcast} esgota as tentativas de
envio. As linhas \texttt{S(FAIL)} indicam falhas percebidas pelo gerador de
sensor ao chamar \texttt{send\_data}. O analisador usa esses registros para
diferenciar perda de amostras por falha de enfileiramento ou transmissão de perda
observada apenas no lado receptor.

As linhas \texttt{WCAN\_BATCH} registram, por identificador CAN, a frequência
média observada, a quantidade média, mínima e máxima de pontos por lote, o
intervalo entre lotes, o tempo médio de despacho e a espera média na fila de
rádio. A instrumentação \texttt{MEASURE} acrescenta as linhas
\texttt{WCAN\_MEASURES}, com tempo decorrido, memória \emph{heap} livre no início
e no fim do ensaio, mínimos de \emph{heap} observados, maior bloco livre, tempo
total estimado de \emph{airtime}, total de pacotes enviados e total de falhas de
envio do sensor.

Essa instrumentação não é neutra. Chamadas a \texttt{esp\_timer\_get\_time()},
seções críticas protegidas por mutex, atualização de contadores e impressão
serial consomem CPU, memória e tempo de execução. Por esse motivo, a campanha
final comparou \emph{broadcast} e \emph{multicast} sob a mesma condição de
instrumentação, ambas com \texttt{MEASURE} habilitado. Os valores coletados devem
ser lidos como resultados do sistema medido, não como limite absoluto de uma
\emph{build} sem coleta. A opção por comprimir intervalos e adiar a maior parte
das impressões para o encerramento reduz a interferência da medição, mas não a
elimina; essa interferência faz parte do custo experimental considerado na
análise.
